% ============ Chapter 4. Valuation rings (原书 pp.71-92) ============
\chapter{Valuation rings}\label{ch:4}

From Hensel's theory of $p$-adic numbers onwards, valuation theory has
been an important tool of number theory and the theory of function fields in
one variable; the main object of study was however the multiplicative
valuations which generalise the usual notion of absolute value of a number.
In contrast, Krull defined and studied valuation rings from a more ring-
theoretic point of view (\cite{Kr3}, 1931). His theory was immediately applied to
algebraic geometry by Zariski. In \S\ref{sec:10} we treat the elementary parts of their
theory. Discrete valuation rings (DVRs) and Dedekind rings, the classical
objects of study, are treated in the following \S\ref{sec:11}, which also includes the
Krull--Akizuki theorem, so that this section contains the theory of one-
dimensional Noetherian rings. In \S\ref{sec:12} we treat Krull rings, which should be
thought of as a natural extension of Dedekind rings; we go as far as a recent
theorem of J.\ Nishimura. This book is mainly concerned with Noetherian
rings, and general valuation rings and Krull rings are the most important
rings outside this category. The present chapter is intended as
complementary to the theory of Noetherian rings, and we have left out
quite a lot on valuation theory. The reader should consult [B6,7],~\cite{ZS} or
other textbooks for more information.

% ============ §10 General valuations (原书 pp.71-78) ============
\section{General valuations}\label{sec:10}\index{valuations, section 10}

An integral domain $R$ is a \emph{valuation ring}\index{valuation ring} if every element $x$ of its
field of fractions $K$ satisfies
\[
  x \notin R \Rightarrow x^{-1} \in R;
\]
if we write $R^{-1}$ for the set of inverses of non-zero elements of $R$ then this
condition can be expressed as $R \cup R^{-1} = K$. We also say that $R$ is a
\emph{valuation ring of $K$}. The case $R = K$ is the \emph{trivial} valuation ring.

If $R$ is a valuation ring then for any two ideals $I, J$ of $R$ either $I \subset J$
or $J \subset I$ must hold; indeed, if $x \in I$ and $x \notin J$ then for any $0 \neq y \in J$ we
have $x/y \notin R$, so that $y/x \in R$ and $y = x(y/x) \in I$, therefore $J \subset I$. Thus the
ideals of $R$ form a totally ordered set. In particular, since $R$ has only one
maximal ideal, $R$ is a local ring. We write $m$ for the maximal ideal of $R$.
Then as one sees easily,
\[
  K - R = \{x \in K^* \mid x^{-1} \in m\},
\]
where we write $K^*$ for the multiplicative group $K - \{0\}$. Thus $R$ is determined by $R$
and $m$.

If $R$ is a valuation ring of a field $K$ then any ring $R'$ with $R \subset R' \subset K$ is
obviously also a valuation ring, and in fact we have the following stronger
statement.

\begin{theorem}\label{thm:10.1}
Let $R \subset R' \subset K$ be as above, let $m$ be the maximal ideal of $R$
and $p$ that of $R'$, and suppose that $R \neq R'$. Then
\begin{enumerate}[(i)]
  \item $p \subset m \subset R \subset R'$ and $p \neq m$.
  \item $p$ is a prime ideal of $R$ and $R' = R_p$.
  \item $R/p$ is a valuation ring of the field $R'/p$.
  \item Given any valuation ring $\bar{S}$ of the field $R'/p$, let $S$ be its inverse image
        in $R'$. Then $S$ is a valuation ring having the same field of fractions $K$ as $R'$.
\end{enumerate}
\end{theorem}

\begin{proof}
(i) If $x \in p$ then $x^{-1} \notin R'$, so $x^{-1} \notin R$ and hence $x \in m$; $x$ is not a unit of $R$,
so that $x \in m$. Also, since $R \neq R'$ we have $p \neq m$.

(ii) We know that $p \subset R$, so that $p = p \cap R$, and this is a prime ideal of $R$.
Since $R - p \subset R' - p = \{\text{units of } R'\}$ we have $R_p \subset R'$, and moreover by
construction, the maximal ideal of $R_p$ is contained in the maximal ideal $p$ of
$R'$. Thus by (i), $R_p = R'$.

(iii) Write $\varphi\colon R' \to R'/p$ for the natural map; then for $x \in R' - p$, if $x \in R$
we have $\varphi(x) \in R/p$, and if $x \notin R$ we have $\varphi(x)^{-1} = \varphi(x^{-1}) \in R/p$, and
therefore $R/p$ is a valuation ring of $R'/p$.

(iv) Note that $p \subset S$ and $S/p = \bar{S}$, so that if $x \in R'$ and $x \notin S$ then $x$ is
a unit of $R'$, and $\varphi(x) \notin \bar{S}$. Thus $\varphi(x^{-1}) = \varphi(x)^{-1} \in \bar{S}$, and hence $x^{-1} \in S$.
If on the other hand $x \in K - R'$ then $x^{-1} \in p \subset S$, so that we have proved
that $S \cup S^{-1} = K$.
\end{proof}

The valuation ring $S$ in (iv) is called the \emph{composite} of $R'$ and $\bar{S}$. According
to (iii), every valuation ring of $K$ contained in $R'$ is obtained as the
composite of $R'$ and a valuation ring of $R'/p$.

Quite generally, we write $m_R$ for the maximal ideal of a local ring $R$. If $R$
and $S$ are local rings with $R \supset S$ and $m_R \cap S = m_S$ we say that $R$ \emph{dominates} $S$,\index{domination}
and write $R \ge S$. If $R \ge S$ and $R \neq S$, we write $R > S$.

\begin{theorem}\label{thm:10.2}
Let $K$ be a field, $A \subset K$ a subring, and $p$ a prime ideal of $A$.
Then there exists a valuation ring $R$ of $K$ satisfying
\[
  R \supset A \quad \text{and} \quad m_R \cap A = p.
\]
\end{theorem}

\begin{proof}
Replacing $A$ by $A_p$ we can assume that $A$ is a local ring with $p = m_A$.
Now write $\mathcal{F}$ for the set of all subrings $B$ of $K$ containing $A$ and such that
$1 \notin pB$. Now $A \in \mathcal{F}$, and if $\mathcal{F}_0 \subset \mathcal{F}$ is a subset totally ordered by inclusion
then the union of all the elements of $\mathcal{F}_0$ is again an element of $\mathcal{F}$, so that, by
Zorn's lemma, $\mathcal{F}$ has an element $R$ which is maximal for inclusion. Since
$pR \neq R$ there is a maximal ideal $m$ of $R$ containing $pR$. Then $R \subset R_m \in \mathcal{F}$, so
that $R = R_m$, and $R$ is local. Also $p \subset m$ and $p$ is a maximal ideal of $A$, so that
$m \cap A = p$. Thus it only remains to prove that $R$ is a valuation ring of $K$. If
$x \in K$ and $x \notin R$ then since $R[x] \notin \mathcal{F}$ we have $1 \in pR[x]$, and there
exists a relation of the form
\[
  1 = a_0 + a_1x + \cdots + a_nx^n \quad \text{with } a_i \in pR.
\]
Since $1 - a_0$ is a unit of $R$ we can modify this to get a relation
\[
  (*) \quad 1 = b_1x + \cdots + b_nx^n \quad \text{with } b_i \in m.
\]
Among all such relations, choose one for which $n$ is as small as possible.
Similarly, if $x^{-1} \notin R$ we can find a relation
\[
  (**) \quad 1 = c_1x^{-1} + \cdots + c_mx^{-m} \quad \text{with } c_i \in m,
\]
and choose one for which $m$ is as small as possible. If $n \ge m$ we multiply
$(**)$ by $b_nx^n$ and subtract from $(*)$, and obtain a relation of the form $(*)$
but with a strictly smaller degree $n$, which is a contradiction; if $n < m$ then
we get the same contradiction on interchanging the roles of $x$ and $x^{-1}$.
Thus if $x \notin R$ we must have $x^{-1} \in R$.
\end{proof}

\begin{theorem}\label{thm:10.3}
A valuation ring is integrally closed.
\end{theorem}

\begin{proof}
Let $R$ be a valuation ring of a field $K$, and let $x \in K$ be integral over $R$,
so that $x^n + a_1x^{n-1} + \cdots + a_n = 0$ with $a_i \in R$. If $x \notin R$ then $x^{-1} \in m_R$, but then
\[
  1 + a_1x^{-1} + \cdots + a_nx^{-n} = 0,
\]
and we get $1 \in m_R$, which is a contradiction. Hence $x \in R$.
\end{proof}

\begin{theorem}\label{thm:10.4}
Let $K$ be a field, $A \subset K$ a subring, and let $B$ be the integral
closure of $A$ in $K$. Then $B$ is the intersection of all the valuation rings of $K$
containing $A$.
\end{theorem}

\begin{proof}
Write $B'$ for the intersection of all valuation rings of $K$ containing
$A$, so that by the previous theorem we have $B' \supset B$. To prove the opposite
inclusion it is enough to show that for any element $x \in K$ which is not inte-
gral over $A$ there is a valuation ring of $K$ containing $A$ but not $x$. Set
$x^{-1} = y$. The ideal $yA[y]$ of $A[y]$ does not contain $1$: for if $1 = a_1y +
a_2y^2 + \cdots + a_ny^n$ with $a_i \in A$ then $x$ would be integral over $A$, contradicting
the assumption. Therefore there is a maximal ideal $p$ of $A[y]$ containing
$yA[y]$, and by Theorem 2 there exists a valuation ring $R$ of $K$ such that
$R \supset A[y]$ and $m_R \cap A[y] = p$. Now $y = x^{-1} \in m_R$, so that $x \notin R$.
\end{proof}

Let $K$ be a field and $A \subset K$ a subring. If a valuation ring $R$ of $K$ contains $A$
we say that $R$ has a \emph{centre} in $A$,\index{centre} and the prime ideal $m_R \cap A$ of $A$ is called the
\emph{centre of $R$ in $A$}. The set of valuation rings of $K$ having a centre in $A$ is called
the \emph{Zariski space}\index{Zariski space} or the \emph{Zariski Riemann surface} of $K$ over $A$, and written
$\Zar(K, A)$. We will treat $\Zar(K, A)$ as a topological space, introducing a
topology as follows.

For $x_1, \ldots, x_n \in K$, set $U(x_1, \ldots, x_n) = \Zar(K, A[x_1, \ldots, x_n])$. Then since
\[
  U(x_1, \ldots, x_n) \cap U(y_1, \ldots, y_m) = U(x_1, \ldots, x_n, y_1, \ldots, y_m),
\]
the collection $\mathcal{B} = \{U(x_1, \ldots, x_n) \mid n \ge 0 \text{ and } x_i \in K\}$ is the basis for
the open sets of a topology on $\Zar(K, A)$. That is, we take as open sets the
subsets of $\Zar(K, A)$ which can be written as a union of elements of $\mathcal{B}$.
As in the case of $\Spec$, this topology is called the \emph{Zariski topology}.

\begin{theorem}\label{thm:10.5}
$\Zar(K, A)$ is quasi-compact.
\end{theorem}

\begin{proof}
It is enough to prove that if $\mathcal{A}$ is a family of closed sets of $\Zar(K, A)$
having the finite intersection property (that is, the intersection of any finite
number of elements of $\mathcal{A}$ is non-empty) then the intersection of all the
elements of $\mathcal{A}$ is non-empty. By Zorn's lemma there exists a maximal family
of closed sets $\mathcal{A}'$ having the finite intersection property and containing $\mathcal{A}$.
Since it is then enough to show that the intersection of all the elements of
$\mathcal{A}'$ is non-empty, we can take $\mathcal{A} = \mathcal{A}'$. Then it is easy to see that $\mathcal{A}$ has the
following properties:
\begin{enumerate}[($\alpha$)]
  \item $F_1, \ldots, F_r \in \mathcal{A} \Rightarrow F_1 \cap \cdots \cap F_r \in \mathcal{A}$,
  \item $Z_1, \ldots, Z_r$ are closed sets and $Z_1 \cup \cdots \cup Z_r \in \mathcal{A} \Rightarrow Z_i \in \mathcal{A}$ for some $i$;
  \item if a closed set $F$ contains an element of $\mathcal{A}$ then $F \in \mathcal{A}$.
\end{enumerate}
For a subset $F \subset \Zar(K, A)$ we write $F^c$ to denote the complement of $F$.
If $F \in \mathcal{A}$ and $F^c = \bigcup_n U_n$ then $F = \bigcap_n U_n^c$, and moreover if $U(x_1, \ldots, x_n)^c =
\bigcup_{i=1}^{n} U(x_i)^c \in \mathcal{A}$ then by ($\alpha$) above one of the $U(x_i)^c$ must belong to $\mathcal{A}$.
Hence the intersection of all the elements of $\mathcal{A}$ is the same thing as the
intersection of the sets of the form $U(x)^c$ belonging to $\mathcal{A}$. Set
\[
  \Gamma = \{y \mid U(y^{-1})^c \in \mathcal{A}\}.
\]
Now since the condition for $R \in \Zar(K, A)$ to belong to $U(y^{-1})^c$ is that
$y \in m_R$, the intersection of all the elements of $\mathcal{A}$ is equal to
\[
  \{R \in \Zar(K, A) \mid m_R \supset \Gamma\}.
\]
Write $I$ for the ideal of $A[\Gamma]$ generated by $\Gamma$. If $1 \notin I$ then by Theorem 2, the
above set is non-empty, as required to prove. But if $1 \in I$ then there is a
finite subset $\{y_1, \ldots, y_n\} \subset \Gamma$ such that $1 \in \sum y_iA[y_1, \ldots, y_n]$; but then
$U(y_1^{-1})^c \cap \cdots \cap U(y_n^{-1})^c = \emptyset$, which contradicts the finite intersection
property of $\mathcal{A}$.
\end{proof}

When $K$ is an algebraic function field over an algebraically closed field $k$
of characteristic $0$ (that is, $K$ is a finitely generated extension of $k$), Zariski
gave a classification of valuation rings of $K$ containing $k$, and using this and
the compactness result above, he succeeded in giving an algebraic proof in
characteristic $0$ of the resolution of singularities of algebraic varieties of
dimension $2$ and $3$. However, Hironaka's general proof of resolution of
singularities in characteristic $0$ in all dimensions was obtained by other
methods, without the use of valuation theory.

As we saw at the beginning of this section, the ideals of $R$ form a totally
ordered set under inclusion. This holds not just for ideals, but for all $R$-
modules contained in $K$. In particular, if we set
\[
  G = \{xR \mid x \in K \text{ and } x \neq 0\},
\]
then $G$ is a totally ordered set under inclusion; we will, however, give $G$ the
opposite order to that given by inclusion. That is, we define $\le$ by
\[
  xR \le yR \iff xR \supset yR.
\]
Moreover, $G$ is an Abelian group with product $(xR) \cdot (yR) = xyR$. In general,
an Abelian group $H$ written additively, together with a total order relation
$\le$ is called an \emph{ordered group}\index{ordered group} if the axiom
\[
  x \ge y, \; z \ge t \Rightarrow x + z \ge y + t
\]
holds. This axiom implies
\[
  \text{(1) } x > 0, \; y > 0 \Rightarrow x + y > 0, \qquad \text{and} \qquad \text{(2) } x > y \Rightarrow -y > -x.
\]
We make an ordered set $H \cup \{\infty\}$ by adding to $H$ an element $\infty$ bigger
than all the elements of $H$, and fix the conventions $\infty + \alpha = \infty$ for $\alpha \in H$
and $\alpha + \infty = \infty$. A map $v\colon K \to H \cup \{\infty\}$ from a field $K$ to $H \cup \{\infty\}$
is called an \emph{additive valuation}\index{additive valuation} or just a \emph{valuation} of $K$ if it satisfies the
conditions
\begin{enumerate}[(1)]
  \item $v(xy) = v(x) + v(y)$;
  \item $v(x + y) \ge \min\{v(x), v(y)\}$;
  \item $v(x) = \infty \iff x = 0$.
\end{enumerate}
If we write $K^*$ for the multiplicative group of $K$ then $v$ defines a
homomorphism $K^* \to H$; the image is a subgroup of $H$, called the \emph{value group} of $v$. We also set
\[
  R_v = \{x \in K \mid v(x) \ge 0\} \quad \text{and} \quad m_v = \{x \in K \mid v(x) > 0\},
\]
obtaining a valuation ring $R_v$ of $K$ with $m_v$ as its maximal ideal, and call $R_v$ the
\emph{valuation ring of $v$}, and $m_v$ the \emph{valuation ideal} of $v$. Conversely, if $R$ is a
valuation ring of $K$, then the group $G = \{xR \mid x \in K^*\}$ described above is an
ordered group, and we obtain an additive valuation of $K$ with value group
$G$ by defining $v\colon K \to G \cup \{\infty\}$ by $v(0) = \infty$ and $v(x) = xR$ for $x \in K^*$ (there
is no real significance in whether or not we rewrite the multiplication in
$G$ additively); the valuation ring of $v$ is $R$. The additive valuation corres-
ponding to a valuation ring $R$ is not quite unique, but if $v$ and $v'$ are
additive valuations of $K$ with value groups $H$ and $H'$ and both having
valuation ring $R$ then there exists an order-isomorphism
$\varphi\colon H \to H'$ such that $v' = \varphi v$ (prove this!). Thus we can think of valuation rings and
additive valuations as being two aspects of the same thing.

We now give some examples of ordered groups:
\begin{enumerate}[(1)]
  \item the additive group of real numbers $\R$ (this is isomorphic to the
        multiplicative group of positive reals), or any subgroup of this;
  \item the group $\Z$ of rational integers;
  \item the direct product $\Z^n$ of $n$ copies of $\Z$, with lexicographical order, that is
        \[
          (a_1, \ldots, a_n) < (b_1, \ldots, b_n) \iff \text{the first non-zero element of}
        \]
        \[
          b_1 - a_1, \ldots, b_n - a_n \text{ is positive}
        \]
\end{enumerate}

An ordered group $G$ is said to be \emph{Archimedean}\index{Archimedean ordered group} if it is order-isomorphic to
a suitable subgroup of $\R$. The name is explained by the following theorem:
the condition in it is known as the Archimedean axiom. (Note that our
usage is completely unrelated to that in number theory, where non-
Archimedean fields are p-adic fields, as opposed to subfields of $\R$ and $\C$ with
the usual `Archimedean' metrics.)

\begin{theorem}\label{thm:10.6}
Let $G$ be an ordered group; then $G$ is Archimedean if and
only if the following condition holds:
\[
  \text{if } a, b \in G \text{ with } a > 0, \text{ there exists a natural number } n \text{ such that } na > b.
\]
\end{theorem}

\begin{proof}
The condition is obviously necessary, and we prove sufficiency. If
$G = \{0\}$ then we can certainly embed $G$ in $\R$. Suppose that $G \neq \{0\}$.
Fix some $0 < \alpha \in G$. For any $\beta \in G$, there is a well-defined largest integer
$n$ such that $n\alpha \le \beta$ (if $\beta \ge 0$ this is clear by assumption; if $\beta < 0$, let $m$ be
the smallest integer such that $-\beta \le m\alpha$, and set $n = -m$). Let this be $n_0$.
Now set $\beta_1 = \beta - n_0\alpha$ and let $n_1$ be the largest integer $n$ such that $n\alpha \le 10\beta_1$;
we have $0 \le n_1 < 10$. Set $\beta_2 = 10\beta_1 - n_1\alpha$ and let $n_2$ be the largest integer
$n$ such that $n\alpha \le 10\beta_2$. Continuing in the same way, we find integers $n_0$,
$n_1, n_2, \ldots$, and set $\varphi(\beta) = \alpha_0$, where $\alpha_0$ is the real number given by the decimal
expression $n_0 + 0.n_1n_2n_3\cdots$. Then it can easily be checked that $\varphi\colon G \to \R$
is order-preserving, in the sense that $\beta < \beta'$ implies $\varphi(\beta) < \varphi(\beta')$.

We also see that $\varphi$ is injective. For this, we only need to observe that
if $\beta < \beta'$ then there exists a natural number $r$ such that $\alpha < 10^r(\beta' - \beta)$;
the details are left to the reader.

Finally we show that $\varphi$ is a group homomorphism. For $\beta \in G$, we write
$n/10^r$ with $n \in \Z$ to denote the number obtained by taking the first $r$
decimal places of $\varphi(\beta)$; the numerator $n$ is determined by the property that
$n\alpha \le 10^r\beta < (n + 1)\alpha$. For $\beta' \in G$, if $n'\alpha \le 10^r\beta' < (n' + 1)\alpha$ then we have
\[
  (n + n')\alpha \le 10^r(\beta + \beta') < (n + n' + 2)\alpha,
\]
and hence
\[
  \varphi(\beta + \beta') - (n + n') \cdot 10^{-r} < 2 \cdot 10^{-r},
\]
so that
\[
  |\varphi(\beta + \beta') - \varphi(\beta) - \varphi(\beta')| \le 4 \cdot 10^{-r},
\]
and since $r$ is arbitrary, $\varphi(\beta + \beta') = \varphi(\beta) + \varphi(\beta')$.
\end{proof}

A non-zero group $G$ order-isomorphic to a subgroup of $\R$ is said to have
\emph{rank 1}.\index{rank (of an ordered group)} The \emph{rational rank} of an ordered group $G$ of rank $1$ is the maximum
number of elements of $G$ (viewed as a subgroup of $\R$) which are linearly
independent over $\Z$. For example, the additive group $G = \Z + \Z\sqrt{2} \subset \R$ is
an ordered group of rank $1$ and rational rank $2$.

\begin{theorem}\label{thm:10.7}
Let $R$ be a valuation ring having value group $G$. Then $G$ has
rank $1 \iff R$ has Krull dimension $1$.
\end{theorem}

\begin{proof}
($\Rightarrow$) Since $G \neq 0$, $R$ is not a field. Suppose that $p$ is a prime ideal of $R$
distinct from $m_R$. Let $\xi \in m_R$ be such that $\xi \notin p$, and set $v(\xi) = x$, where $v$ is the
additive valuation corresponding to $R$. Suppose that $0 \neq \eta \in p$, and set $\psi =
v(\eta)$; then $\psi \in G$ and $x > 0$, so that $nx > \psi$ for some sufficiently large natural
number $n$. This means that $\xi^n/\eta \in R$, so that $\xi^n \in \eta R \subset p$; then since $p$ is prime
we have $\xi \in p$, which is a contradiction. Therefore $p = (0)$. The only prime
ideals of $R$ are $m_R$ and $(0)$, which means $\dim R = 1$.

($\Leftarrow$) If $0 \neq \eta \in m_R$ then $m_R$ is the unique prime ideal containing $\eta R$, and
hence $\sqrt{\eta R} = m_R$. Thus for any $\xi \in m_R$ there exists a natural number $n$ such
that $\xi^n \in \eta R$. From this one sees easily that $G$ satisfies the Archimedean
axiom.
\end{proof}

% ============ 习题 ============
\section*{Exercises to \S\ref{sec:10}}
\emph{Prove the following propositions.}

\begin{exercise}\label{ex:10.1-ex}
In a valuation ring any finitely generated ideal is principal.
\end{exercise}

\begin{exercise}\label{ex:10.2-ex}
If $R$ is a valuation ring then an $R$-module $M$ is flat if and only if it is torsion-
free (that is, $a \neq 0, x \neq 0 \Rightarrow ax \neq 0$ for $a \in R, x \in M$).
\end{exercise}

\begin{exercise}\label{ex:10.3-ex}
In Theorem~\ref{thm:10.4}, if $A$ is a local ring then $B$ is the intersection of the
valuation rings of $K$ dominating $A$.
\end{exercise}

\begin{exercise}\label{ex:10.4-ex}
If $R$ is a valuation ring of Krull dimension $\ge 2$ then the formal power series
ring $R[[X]]$ is not integrally closed ([B5], Ex. 27, p.~76 and Seidenberg~\cite{Sd1}).
\end{exercise}

\begin{exercise}\label{ex:10.5-ex}
If $R$ is a valuation ring of Krull dimension $1$ and $K$ its field of fractions then
there do not exist any rings intermediate between $R$ and $K$. In other words
$R$ is maximal among proper subrings of $K$. Conversely if a ring $R$, not a
field, is a maximal proper subring of a field $K$ then $R$ is a valuation ring of
Krull dimension $1$.
\end{exercise}

\begin{exercise}\label{ex:10.6-ex}
If $v$ is an additive valuation of a field $K$, and if $\alpha, \beta \in K$ are such that
$v(\alpha) \neq v(\beta)$ then $v(\alpha + \beta) = \min\{v(\alpha), v(\beta)\}$.
\end{exercise}

\begin{exercise}\label{ex:10.7-ex}
If $v$ is an additive valuation of a field $K$ and $\alpha_1, \ldots, \alpha_n \in K$ are such that
$\alpha_1 + \cdots + \alpha_n = 0$ then there exist two indices $i, j$ such that $i \neq j$ and
$v(\alpha_i) = v(\alpha_j)$.
\end{exercise}

\begin{exercise}\label{ex:10.8-ex}
Let $K \subset L$ be algebraic field extension of degree $[L:K] = n$, and let $S$ be
a valuation ring of $L$; set $R = S \cap K$. Write $k, k'$ for the residue fields of $S$
and $R$, and set $[k:k'] = f$. Now let $G$ be the value group of $S$, and let $G'$ be
the image of $K^*$ under the valuation map $L^* \to G$; set $[G:G'] = e$. Then
$ef \le n$. (The numbers $f$ and $e$ are called the \emph{degree} and the \emph{ramification}
\emph{index} of the valuation ring extension $S/R$.)
\end{exercise}

\begin{exercise}\label{ex:10.9-ex}
Let $L, K, S$ and $R$ be as in the previous question, and let $S_1 \neq S$ be a
valuation ring of $L$ such that $S_1 \cap K = R$. Then neither of $S$ or $S_1$ contains
the other.
\end{exercise}

\begin{exercise}\label{ex:10.10-ex}
Let $A$ be an integral domain with field of fractions $K$, and let $H$ be an
ordered group. If a map $v\colon A \to H \cup \{\infty\}$ satisfies conditions (1), (2) and (3)
of an additive valuation (on elements of $A$), then $v$ can be extended
uniquely to an additive valuation $K \to H \cup \{\infty\}$.
\end{exercise}

\begin{exercise}\label{ex:10.11-ex}
Let $k$ be a field, $X$ and $Y$ indeterminates, and suppose that $\alpha$ is a positive
irrational number. Then the map $v\colon k[X, Y] \to \R \cup \{\infty\}$ defined by
taking $\sum c_{nm}X^nY^m$ (with $c_{nm} \in k$) into $v(\sum c_{nm}X^nY^m) = \min\{n +
m\alpha \mid c_{nm} \neq 0\}$ determines a valuation of $k(X, Y)$ with value group
$\Z + \Z\alpha$.
\end{exercise}

% ============ §11 DVRs and Dedekind rings (原书 pp.78-86) ============
\section{DVRs and Dedekind rings}\label{sec:11}\index{discrete valuation rings, section 11}\index{Dedekind rings, section 11}

A valuation ring whose value group is isomorphic to $\Z$ is called
a \emph{discrete valuation ring} (DVR).\index{discrete valuation ring} Discrete refers to the fact that the value
group is a discrete subgroup of $\R$, and has nothing to do with the
$m$-adic topology of the local ring being discrete.

\begin{theorem}\label{thm:11.1}
Let $R$ be a valuation ring. Then the following conditions are
equivalent.
\begin{enumerate}[(1)]
  \item $R$ is a DVR;
  \item $R$ is a principal ideal domain;
  \item $R$ is Noetherian.
\end{enumerate}
\end{theorem}

\begin{proof}
Let $K$ be the field of fractions of $R$ and $m$ its maximal ideal.

$(1) \Rightarrow (2)$ Let $v_R$ the additive valuation of $R$ having value group $\Z$; this is
called the \emph{normalised additive valuation} corresponding to $R$. There exists
$t \in R$ such that $v_R(t) = 1$. For $0 \neq x \in m$, the valuation $v_R(x)$ is a positive
integer, say $v_R(x) = n$; then $v_R(x/t^n) = 0$, so that we can write $x = t^nu$ with $u$ a
unit of $R$. In particular $m = tR$. Let $I \neq (0)$ be any ideal of $R$; then
$\{v_R(a) \mid 0 \neq a \in I\}$ is a set of non-negative integers, and so has a smallest
element, say $n$. If $n = 0$ then $I$ contains a unit of $R$, so that $I = R$. If $n > 0$ then
there exists an $x \in I$ such that $v_R(x) = n$; then $I = xR = t^nR$. Therefore $R$ is a
principal ideal domain, and moreover every non-zero ideal of $R$ is a power
of $m = tR$.

$(2) \Rightarrow (3)$ is obvious.

$(3) \Rightarrow (2)$ In general, given any two ideals of a valuation ring, one contains
the other, so that any finitely generated ideal $a_1R + \cdots + a_nR$ is equal to one
of the $a_iR$, and therefore principal. Hence, if $R$ is Noetherian every ideal of $R$
is principal.

$(2) \Rightarrow (1)$ We can write $m = xR$ for some $x$. Now if we set $I = \bigcap_{\nu \ge 1} x^\nu R$
then this is also a principal ideal, so that we can write $I = yR$. If we set
$y = xz$, then from $y \in x^\nu R$ we get $z \in x^{\nu-1}R$, and since this holds for every
$\nu$, we have $z \in I$, hence we can write $z = yu$. Since $y = xz = xyu$, we have $y(1
- xu) = 0$, but then since $x \in m$ we must have $y = 0$, and therefore $I = (0)$.
Because of this, for every non-zero element $a \in R$, there is a well-defined
integer $\nu \ge 0$ such that $a \in x^\nu R$ but $a \notin x^{\nu+1}R$; we then set $v(a) = \nu$. It is not
difficult to see that if $a, b, c, d \in R - \{0\}$ satisfy $a/b = c/d$ then
\[
  v(a) - v(b) = v(c) - v(d);
\]
therefore setting $v(\xi) = v(a) - v(b)$ for $\xi = a/b \in K^*$ gives a map $v\colon K^* \to \Z$
which can easily be seen to be an additive valuation of $K$ whose valuation
ring is $R$. The value group of $v$ is clearly $\Z$, so that $R$ is a DVR.
\end{proof}

If $R$ is a DVR with maximal ideal $m$ then an element $t \in R$ such that
$m = tR$ is called a \emph{uniformising element} of $R$.\index{uniformising element}

\begin{remark}
A valuation ring $S$ whose maximal ideal $m_S$ is principal does not
have to be a DVR. To obtain a counter-example, let $K$ be a field, and $R$
a DVR of $K$; set $k = R/m_R$, and suppose that $\bar{R}$ is a DVR of $k$. Now let
$S$ be the composite of $R$ and $\bar{R}$. Let $f$ be a uniformising element of $R$, and
$g \in S$ be any element mapping to a uniformising element $\bar{g}$ of $\bar{R}$. Then
$\eta = fR \subset m_R \subset S \subset R$, and $m_S/m_R = \bar{g}\bar{R} = \sqrt{\bar{g}S/\bar{m}_R}$, and so
\[
  m_S = m_R + gS.
\]
On the other hand $g^{-1} \in R$, so that for any $h \in m_R$ we have $h/g \in m_R \subset S$, and
hence $m_R \subset gS$, so that
\[
  m_S = gS.
\]
However, $m_R = fR$ is not finitely generated as an ideal of $S$, being generated
by $f, fg^{-1}, fg^{-2}, \ldots$. The value group of $S$ is $\Z^2$, with the valuation
$v\colon K^* \to \Z^2$ given by
\[
  v(x) = (n, m), \quad \text{where } n = v_R(x) \text{ and } m = v_{\bar{R}}(\bar{\varphi}(xf^{-n})),
\]
where $\varphi\colon R \to R/m_R = k$ is the natural homomorphism.
\end{remark}

The previous theorem gives a characterisation of DVRs among valuation
rings; now we consider characterisations among all rings.

\begin{theorem}\label{thm:11.2}
Let $R$ be a ring; then the following conditions are
equivalent:
\begin{enumerate}[(1)]
  \item $R$ is a DVR;
  \item $R$ is a local principal ideal domain, and not a field;
  \item $R$ is a Noetherian local ring, $\dim R > 0$ and the maximal ideal $m_R$ is
        principal;
  \item $R$ is a one-dimensional normal Noetherian local ring.
\end{enumerate}
\end{theorem}

\begin{proof}
We saw $(1) \Rightarrow (2)$ in the previous theorem; $(2) \Rightarrow (3)$ is obvious.

$(3) \Rightarrow (1)$ Let $xR$ be the maximal ideal of $R$. If $x$ were nilpotent then we
would have $\dim R = 0$, and hence $x^\nu \neq 0$ for all $\nu$. By the Krull intersection
theorem (Theorem~\ref{thm:8.10}, (i)) we have $\bigcap_{\nu \ge 1} x^\nu R = (0)$, so that for $0 \neq y \in R$
there is a well-determined $\nu$ such that $y \in x^\nu R$ and $y \notin x^{\nu+1}R$. If $y = x^\nu u$, then
since $u \notin xR$ it must be a unit. Similarly, for $0 \neq z \in R$ we have $z = x^\mu v$, with $v$ a
unit. Therefore $yz = x^{\mu+\nu}uv \neq 0$, and so $R$ is an integral domain. Finally,
any element $\xi$ of the fraction field of $R$ can be written $\xi = x^\nu u$, with $u$ a unit of
$R$ and $\nu \in \Z$, and it is easy to see that setting $v(\xi) = \nu$ defines an additive
valuation of the field of fractions of $R$ whose valuation ring is $R$.

$(1) \Rightarrow (4)$ In a DVR the only ideals are $(0)$ and the powers of the maximal
ideal, so that the only prime ideals of $R$ are $(0)$ and $m_R$, and hence
$\dim R = 1$. By the previous theorem $R$ is Noetherian, and it is normal
because it is a valuation ring.

$(4) \Rightarrow (3)$ By assumption $R$ is an integral domain. Write $K$ for the field of
fractions and $m$ for the maximal ideal of $R$. Then $m \neq 0$, so that by Theorem~\ref{thm:8.10}, $m \neq m^2$; choose some $x \in m - m^2$. Since $\dim R = 1$ the only prime
ideals of $R$ are $(0)$ and $m$, so that $m$ must be a prime divisor of $xR$, and there
exists $y \in R$ such that $xR:y = m$. Set $a = yx^{-1}$; then $a \notin R$, but $am \subset R$. Now
we set $m^{-1} = \{\xi \in K \mid \xi m \subset R\}$, so that $R \subset m^{-1}$, and $R \neq m^{-1}$ since $a \in m^{-1}$.
Consider the ideal $m^{-1}m$ of $R$; since $R \subset m^{-1}$ we have $m \subset m^{-1}m$. If we
had $m = m^{-1}m$ then we would get $am \subset m$, and then $a$ would be integral
over $R$ by Theorem~\ref{thm:2.1}, so that $a \in R$, which is a contradiction. Hence we
must have $m^{-1}m = R$. Moreover, $xm^{-1} \subset R$ is an ideal, and if $xm^{-1} \subset m$
then we would have $xR = xm^{-1}m \subset m^2$, contradicting $x \notin m^2$. Therefore
$xm^{-1} = R$, and hence $xR = xm^{-1}m = m$, so that $m$ is principal.
\end{proof}

Quite generally, if $R$ is an integral domain and $K$ its fields of fractions, we
say that an $R$-submodule $I$ of $K$ is a \emph{fractional ideal}\index{fractional ideal} of $R$ if $I \neq 0$, and
there exists a non-zero element $\alpha \in R$ such that $\alpha I \subset R$ (see Ex.~\ref{ex:3.4-ex}). As
an $R$-module we have $I \cong \alpha I$, so that if $R$ is a Noetherian integral domain
then any fractional ideal is finitely generated. For $I$ a fractional ideal we
set $I^{-1} = \{\alpha \in K \mid \alpha I \subset R\}$; we say that $I$ is \emph{invertible}\index{invertible ideal} if $I^{-1}I = R$.

\begin{theorem}\label{thm:11.3}
Let $R$ be an integral domain and $I$ a fractional ideal of $R$.
Then the following conditions are equivalent:
\begin{enumerate}[(1)]
  \item $I$ is invertible;
  \item $I$ is a projective $R$-module;
  \item $I$ is finitely generated, and for every maximal ideal $P$ of $R$, the
        fractional ideal $I_P = IR_P$ of $R_P$ is principal.
\end{enumerate}
\end{theorem}

\begin{proof}
$(1) \Rightarrow (2)$ If $I^{-1}I = R$ then there exist $a_i \in I$ and $b_i \in I^{-1}$ such that
$\sum_i a_ib_i = 1$. Then $a_1, \ldots, a_n$ generate $I$, since for any $x \in I$ we have $\sum (xb_i)a_i =
x$, and $xb_i \in R$. Let $F = Re_1 + \cdots + Re_n$ be the free $R$-module with basis
$e_1, \ldots, e_n$; we define the $R$-linear map $\varphi\colon F \to I$ by $\varphi(e_i) = a_i$, so that $\varphi$ is
surjective. Then we define $\psi\colon I \to F$ by writing $\psi_i\colon I \to R$ for the map
$\psi_i(x) = b_ix$, and setting $\psi(x) = \sum \psi_i(x)e_i$. We then have $\varphi\psi(x) = x$, so that $\varphi$
splits, and $I$ is isomorphic to a direct summand of the free module $F$, and
therefore projective.

$(2) \Rightarrow (1)$ Every $R$-linear map from $I$ to $R$ is given by multiplication by
some element of $K$ (prove this!). If we let $\varphi\colon F \to I$ be a surjective map from
a free module $F = \bigoplus Re_i$, by assumption there exists a splitting $\psi\colon I \to F$
such that $\varphi\psi = 1$. Write $\psi(x) = \sum \lambda_i(x)e_i$ for $x \in I$; then by what we have said,
each $\lambda_i$ determines a $b_i \in K$ such that $\lambda_i(x) = b_ix$, and since for each $x$ there
are only finitely many $i$ such that $\lambda_i(x) \neq 0$, we have $b_i = 0$ for all but finitely
many $i$. Letting $b_1, \ldots, b_n$ be the non-zero ones, we have $\sum a_ib_ix = x$ for all
$x \in I$, where $a_i = \varphi(e_i)$. Thus $\sum_i a_ib_i = 1$. Moreover, since $\lambda_iI = \lambda_i \subset R$ we
have $b_i \in I^{-1}$, and therefore $I^{-1}I = R$.

$(1) \Rightarrow (3)$ As we have already seen, $I$ is finitely generated. Now if $\sum a_ib_i = 1$
and $P$ is any prime ideal then at least one of $a_ib_i$ must be a unit of $R_P$,
and $I_P = a_iR_P$. Hence $I_P$ is a principal fractional ideal.

$(3) \Rightarrow (1)$ If $I$ is finitely generated then $(I^{-1})_P = (I_P)^{-1}$. Indeed, the
inclusion $\subset$ holds for any ideal; for $\supset$, if $I = a_1R + \cdots + a_nR$ and $x \in (I_P)^{-1}$
then $xa_i \in R_P$, so there exist $c_i \in R - P$ such that $xa_ic_i \in R$, so that setting
$c = c_1 \cdots c_n$ we have $(cx)a_i \in R$ for all $i$, which gives $cx \in I^{-1}$ and $x \in (I^{-1})_P$.
From the fact that $I_P$ is principal, we get $I_P \cdot (I_P)^{-1} = R_P$. Now if $II^{-1} \neq R$
then we can take a maximal ideal $P$ such that $II^{-1} \subset P$, and then
$I_P \cdot (I_P)^{-1} = I_P(I^{-1})_P \subset PR_P$, which is a contradiction. Thus we must have
$II^{-1} = R$.
\end{proof}

\begin{theorem}\label{thm:11.4}
Let $R$ be a Noetherian integral domain, and $P$ a non-zero
prime ideal of $R$. If $P$ is invertible then $\htop P = 1$ and $R_P$ is a DVR.
\end{theorem}

\begin{proof}
If $P$ is invertible the maximal ideal $PR_P$ of $R_P$ is principal, and
condition (3) of Theorem 2 is satisfied; thus $R_P$ is a DVR, and so $\dim R_P = 1$.
\end{proof}

\begin{theorem}\label{thm:11.5}
Let $R$ be a normal Noetherian domain. Then we have
\begin{enumerate}[(i)]
  \item all the prime divisors of a non-zero principal ideal have height $1$;
  \item $R = \bigcap_{\htop p = 1} R_p$.
\end{enumerate}
\end{theorem}

\begin{proof}
(i) Suppose $0 \neq a \in R$ and that $P$ is one of the prime divisors of $aR$;
then there exists an element $b \in R$ such that $aR:b = P$. We set $PR_P = m$, and
then $aR_P:b = m$, so that $ba^{-1} \notin R_P$.
If $ba^{-1}m \subset m$ then by the determinant trick $ba^{-1}$ is integral over $R_P$, which
contradicts the fact that $R_P$ is integrally closed. Thus $ba^{-1}m = R_P$, so that $m^{-1}m = R_P$, and
then by the previous theorem we get $\htop m = \htop P = 1$.

(ii) It is sufficient to prove that for $a, b \in R$ with $a \neq 0$, $b \in aR_P$ for every
height $1$ prime $P \in \Spec R$ implies $b \in aR$. Let $P_1, \ldots, P_r$ be the prime
divisors of $aR$, and let $aR = q_1 \cap \cdots \cap q_r$ be a primary decomposition of
$aR$, where $q_i$ is a $P_i$-primary ideal for each $i$. Then since $\htop P_i = 1$, we have
$b \in aR_{P_i} \cap R = q_i$ for $i = 1, \ldots, n$, and therefore $b \in \bigcap q_i = aR$.
\end{proof}

\begin{corollary*}
Let $R$ be a Noetherian domain. The following two conditions
are necessary and sufficient for $R$ to be normal:
\begin{enumerate}[(a)]
  \item for $P$ a height $1$ prime ideal, $R_P$ is a DVR;
  \item all the prime divisors of a non-zero principal ideal of $R$ have height $1$.
\end{enumerate}
\end{corollary*}

\begin{proof}
We have already seen necessity. For sufficiency, note that the proof
of (ii) above shows that (b) implies $R = \bigcap_{\htop p = 1} R_p$. Then by (a) each $R_p$
is normal, so that $R$ is normal.
\end{proof}

\begin{definition}
An integral domain for which every non-zero ideal is invertible
is called a \emph{Dedekind ring}\index{Dedekind ring} (sometimes \emph{Dedekind domain}).
\end{definition}

\begin{theorem}\label{thm:11.6}
For an integral domain $R$ the following conditions are
equivalent:
\begin{enumerate}[(1)]
  \item $R$ is a Dedekind ring;
  \item $R$ is either a field or a one-dimensional Noetherian normal domain;
  \item every non-zero ideal of $R$ can be written as a product of a finite
        number of prime ideals.
\end{enumerate}
Moreover, the factorisation into primes in (3) is unique.
\end{theorem}

\begin{proof}
$(1) \Rightarrow (2)$ Every non-zero ideal is invertible, and therefore finitely
generated, so that $R$ is Noetherian. Let $P$ be a non-zero prime ideal of $R$;
then by Theorem 4, the local ring $R_P$ is a DVR and $\htop P = 1$, and therefore
either $R$ is a field or $\dim R = 1$. Also, by Theorem~\ref{thm:4.7} we know that $R$ is the
intersection of the $R_P$ as $P$ runs through all the maximal ideals of $R$, but
since each $R_P$ is a DVR it follows that $R$ is normal.

$(2) \Rightarrow (1)$ If $R$ is a field there is no problem. If $R$ is not a field then for every
maximal ideal $P$ of $R$ the local ring $R_P$ is a one-dimensional Noetherian
local ring and is normal, so that by Theorem 2 it is a principal ideal ring.
Thus by Theorem 3, $R$ is a Dedekind ring.

$(1) \Rightarrow (3)$ Let $I$ be a non-zero ideal. If $I = R$ then we can view it as the
product of zero ideals; if $I$ is itself maximal then it is the product of just one
prime ideal. We have already seen that $R$ is Noetherian, so that we can use
descending induction on $l(I)$, that is assume that $I \neq R$ and that every ideal
strictly bigger than $I$ is a product of prime ideals. If $I \neq R$ then there is a
maximal ideal $P$ containing $I$, and $I \subset IP^{-1} \subset R$. If $IP^{-1} = I$ then using
$P^{-1}P = R$ we would have $I = IP$, and by NAK this would lead to a
contradiction. Hence $IP^{-1} \neq I$, so that by induction we can write $IP^{-1} =
Q_1 \cdots Q_s$, with $Q_i \in \Spec R$. Multiplying both sides by $P$ gives $I = Q_1 \cdots Q_sP$.

The proof of $(3) \Rightarrow (1)$ is a little harder, and we break it up into four steps.

\textbf{Step 1.} In general, any non-zero principal ideal $aR$ of an integral domain $R$
is obviously invertible. Moreover, suppose that $I$ and $J$ are non-zero
fractional ideals and $B = IJ$; then obviously $I$ and $J$ invertible implies $B$
invertible, but the converse also holds. To see this, from $I^{-1}J^{-1}B \subset R$ we
get $I^{-1}J^{-1} \subset B^{-1}$, and also from $B^{-1}IJ \subset R$ we get $B^{-1}I \subset J^{-1}$ and
$B^{-1}J \subset I^{-1}$; now if $B$ is invertible then multiplying the last two inclusions
together we get $B^{-1} = B^{-1}B^{-1}IJ \subset I^{-1}J^{-1}$, and hence $B^{-1} = I^{-1}J^{-1}$.
Therefore
\[
  R = BB^{-1} = IJI^{-1}J^{-1} = (II^{-1})(JJ^{-1}),
\]
and we must have $II^{-1} = JJ^{-1} = R$.

\textbf{Step 2.} Let $P$ be a non-zero prime ideal. Let us prove that if $I$ is an ideal
strictly bigger than $P$ then $IP = P$. For this it is sufficient to show that if
$I = P + aR$ with $a \notin P$ then $P \subset IP$. Consider expressions of $P^2$ and $a^2R + P$
as product of prime ideals, $P^2 = P_1 \cdots P_r$ and $a^2R + P = Q_1 \cdots Q_s$. Then $P_i$
and $Q_j$ are prime ideals containing $P^2$, and so are prime ideals strictly bigger than $P$,
We now set $\bar{R} = R/P$, and write $-$ to denote the image in $\bar{R}$ of
elements or ideals of $R$. Then we have
\[
  (*) \quad \bar{P}_1 \cdots \bar{P}_r = \bar{a}^2\bar{R} = \bar{Q}_1 \cdots \bar{Q}_s,
\]
and applying Step 1 to the domain $\bar{R}$ we find that $\bar{P}_i$ and $\bar{Q}_j$ are all invertible,
and are prime ideals of $\bar{R}$. We can suppose that $\bar{P}_1$ is a minimal element of
the set $\{\bar{P}_1, \ldots, \bar{P}_r\}$. Moreover, at least one of $\bar{Q}_1, \ldots, \bar{Q}_s$ is contained
in $\bar{P}_1$, so that we can assume that $\bar{Q}_1 \subset \bar{P}_1$, and, on the other hand, since
$\bar{Q}_1$ is also a prime ideal and $\bar{P}_1 \cdots \bar{P}_r \subset \bar{Q}_1$ we must have $\bar{Q}_1 \supset \bar{P}_i$
for some $i$. Then $\bar{P}_i \subset \bar{Q}_1 \subset \bar{P}_1$, and by the minimality of $\bar{P}_1$ we have $\bar{P}_i =
\bar{P}_1 = \bar{Q}_1$. Multiplying through both sides of $(*)$ by $\bar{P}_1^{-1}$ gives
\[
  \bar{P}_2 \cdots \bar{P}_r = \bar{Q}_2 \cdots \bar{Q}_s;
\]
proceeding in the same way, we see that $r = s$, and that after reordering the
$\bar{Q}_j$ we can assume that $\bar{P}_i = \bar{Q}_i$ for $i = 1, \ldots, r$. From this we get $P_i = Q_i$, and
$a^2R + P = (P + aR)^2 = P^2 + aP + a^2R$. Thus any element $x \in P$ can be
written
\[
  x = y + az + a^2t \quad \text{with } y \in P^2, z \in P \text{ and } t \in R.
\]
Since $a \notin P$ we must have $t \in P$, and then as required we have
$P \subset P^2 + aP = (P + aR)P$.

\textbf{Step 3.} Let $b \in R$ be a non-zero element; then in the factorisation
$bR = P_1 \cdots P_r$, every $P_i$ is a maximal ideal of $R$. Indeed, if $I$ is any ideal
strictly greater than $P_i$ then $IP_i = P_i$, and by Step 1 $P_i$ is invertible, so that
$I = R$.

\textbf{Step 4.} Let $P$ be a prime ideal of $R$, and $0 \neq a \in P$. If $aR = P_1 \cdots P_r$ with
$P_i \in \Spec R$ then $P$ must contain one of the $P_i$, but from Step 3 we know that
$P_i$ is maximal, so that $P = P_i$. We deduce that $P$ is a maximal ideal and is
invertible. If every non-zero prime ideal is invertible then any non-zero ideal
of $R$ is invertible, since it can be written as a product of primes. This
completes the proof of $(3) \Rightarrow (1)$.

Finally, if (1), (2) and (3) hold, then as we have seen in Step 2 above, the
uniqueness of factorisation into primes is a consequence of the fact that
the prime ideals of $R$ are invertible.
\end{proof}

\begin{theorem}[the Krull--Akizuki theorem]\label{thm:11.7}
Let $A$ be a one-dimensional
Noetherian integral domain with field of fractions $K$, let $L$ be a finite
algebraic extension field of $K$, and $B$ a ring with $A \subset B \subset L$; then $B$ is a
Noetherian ring of dimension at most $1$, and if $J$ is a non-zero ideal of $B$ then
$B/J$ is an $A$-module of finite length.
\end{theorem}

\begin{proof}
We follow the method of proof of Akizuki [1] in the linear algebra
formulation of [BS]. First of all we prove the following lemma.

\begin{lemma}\label{lem:11.1}
Let $A$ and $K$ be as in the theorem, and let $M$ be a torsion-free
$A$-module (see Ex.~\ref{ex:10.2-ex}) of rank $r < \infty$. Then for $0 \neq a \in A$ we have
\[
  l(M/aM) \le r \cdot l(A/aA).
\]
\end{lemma}

\begin{remark}
The \emph{rank} of a module $M$ over an integral domain $A$ is the maximal
number of elements of $M$ linearly independent over $A$; this is equal to the
dimension of the $K$-vector space $M \otimes_A K$.
\end{remark}

\begin{proof}[Proof of the lemma]
First we assume that $M$ is finitely generated. Choose
elements $\xi_1, \ldots, \xi_r \in M$ linearly independent over $A$ and set $E = \sum A\xi_i$; then
for any $\eta \in M$ there exists $t \in A$ with $t \neq 0$ such that $t\eta \in E$. If we set $C = M/E$
then from the assumption on $M$ we see that $C$ is also finitely generated, so
that $tC = 0$ for suitable $0 \neq t \in A$. Applying Theorem~\ref{thm:6.4} to $C$, we can find
$C = C_0 \supset C_1 \supset \cdots \supset C_m = 0$, such that $C_{i-1}/C_i \cong A/p_i$, with $p_i \in \Spec A$.
Now $t \in p_i$, and since $A$ is one-dimensional each $p_i$ is maximal, so that
$l(C) = m < \infty$. If $0 \neq a \in A$ then the exact sequence
\[
  E/a^nE \to M/a^nM \to C/a^nC \to 0
\]
gives
\[
  (*) \quad l(M/a^nM) \le l(E/a^nE) + l(C) \quad \text{for all } n > 0.
\]
Now $E$ and $M$ are both torsion-free $A$-modules, and one sees easily that
$aM/a^{i+1}M \cong M/aM$, and similarly for $E$. Hence $(*)$ can be written
$n \cdot l(M/aM) \le n \cdot l(E/aE) + l(C)$ for all $n > 0$, which gives $l(M/aM) \le l(E/aE)$.
Since $E \cong A^r$ we have $l(E/aE) = r \cdot l(A/aA)$. This completes the proof in the
case that $M$ is finitely generated. If $M$ is not finitely generated, take a finitely
generated submodule $\bar{M}_0 = A\bar{\xi}_1 + \cdots + A\bar{\xi}_r$ of $\bar{M} = M/aM$. Then choos-
ing an inverse image $\omega_i$ in $M$ for each $\bar{\xi}_i$, and setting $M_0 = \sum A\omega_i$, we get
\[
  l(\bar{M}_0) = l(M_0/M_0 \cap aM) \le l(M_0/aM_0) \le r \cdot l(A/aA).
\]
The right-hand side is now independent of $\bar{M}_0$, so that $\bar{M}$ is in fact finitely
generated, and $l(\bar{M}) \le r \cdot l(A/aA)$.
\end{proof}

We return to the proof of the theorem. We can replace $L$ by the field of
fractions of $B$. Set $[L:K] = r$; then $B$ is a torsion-free $A$-module of rank $r$.
Hence by the lemma, for any $0 \neq a \in A$ we have $l_A(B/aB) < \infty$. Now if $J \neq 0$ is
an ideal of $B$ and $0 \neq b \in J$ then since $b$ is algebraic over $A$ it satisfies a
relation
\[
  a_0b^m + a_1b^{m-1} + \cdots + a_{m-1}b + a_m = 0 \quad \text{with } a_i \in A.
\]
$B$ is an integral domain, so that we can assume $a_0 \neq 0$. Then $0 \neq a_0 \in J \cap A$
and so
\[
  l_A(B/J) \le l_A(B/a_0B) < \infty.
\]
Moreover, one sees from $l_A(J/a_0B) \le l_A(B/a_0B) < \infty$ that $J/a_0B$
is a finite $B$-module; hence, $J$ itself is a finite $B$-module, and therefore $B$ is
Noetherian. If $P$ is a non-zero prime ideal of $B$ then $B/P$ is an Artinian
ring and an integral domain, and therefore a field. Thus $P$ is maximal and
$\dim B = 1$.
\end{proof}

\begin{corollary*}
Let $A$ be a one-dimensional Noetherian integral domain, $K$ its
field of fractions, and $L$ a finite algebraic extension field of $K$; write $B$ for the
integral closure of $A$ in $L$. Then $B$ is a Dedekind ring, and for any maximal
ideal $P$ of $A$ there are just a finite number of primes of $B$ lying over $P$.
\end{corollary*}

\begin{proof}
By the theorem $B$ is a one-dimensional Noetherian integral domain
and is normal by construction; hence it is a Dedekind ring. It is easy to see
that if we factorise $PB$ as a product $PB = Q_1^{e_1} \cdots Q_r^{e_r}$ of a finite number of
prime ideals, then $Q_1, \ldots, Q_r$ are all the prime ideals of $B$ lying over $P$.
\end{proof}

% ============ 习题 ============
\section*{Exercises to \S\ref{sec:11}}
\emph{Prove the following propositions.}

\begin{exercise}\label{ex:11.1-ex}
Let $A$ be a DVR, $K$ its field of fractions, and $\bar{K}$ an algebraic closure of $K$;
then any valuation ring of $\bar{K}$ dominating $A$ is a one-dimensional non-
discrete valuation ring.
\end{exercise}

\begin{exercise}\label{ex:11.2-ex}
Let $A$ be a DVR, $K$ its field of fractions, and $L$ a finite extension field of $K$;
then a valuation ring of $L$ dominating $A$ is a DVR.
\end{exercise}

\begin{exercise}\label{ex:11.3-ex}
Let $A$ be a DVR and $m$ its maximal ideal; then the $m$-adic completion $\hat{A}$ of
$A$ is again a DVR.
\end{exercise}

\begin{exercise}\label{ex:11.4-ex}
Let $v\colon K \to \R \cup \{\infty\}$ be an Archimedean additive valuation of a field $K$,
and let $c$ be a real number with $0 < c < 1$. For $\alpha, \beta \in K$, set $d(\alpha, \beta) = c^{v(\alpha-\beta)}$;
then $d$ satisfies the axioms for a metric on $K$ (that is $d(\alpha, \beta) \ge 0$, $d(\alpha, \beta)
= 0 \iff \alpha = \beta$, $d(\alpha, \beta) = d(\beta, \alpha)$ and $d(\alpha, \gamma) \le d(\alpha, \beta) + d(\beta, \gamma)$), and the to-
pology of $K$ defined by $d$ does not depend on the choice of $c$. Let $R$ be the
valuation ring of $v$ and $m$ its valuation ideal; if $R$ is a DVR then the
topology determined by $d$ restricts to the $m$-adic topology on $R$.
\end{exercise}

\begin{exercise}\label{ex:11.5-ex}
Any ideal in a Dedekind ring can be generated by at most two elements.
\end{exercise}

\begin{exercise}\label{ex:11.6-ex}
Let $A$ be the integral closure of $\Z$ in $\Q(\sqrt{10})$; then $A$ is a Dedekind ring but
not a principal ideal ring.
\end{exercise}

\begin{exercise}\label{ex:11.7-ex}
If a Dedekind ring $A$ is semilocal then it is a principal ideal ring.
\end{exercise}

\begin{exercise}\label{ex:11.8-ex}
A module over a Dedekind ring is flat if and only if it is torsion-free.
\end{exercise}

\begin{exercise}\label{ex:11.9-ex}
Let $A$ be an integral domain (not necessarily Noetherian). The following
two conditions are equivalent:
\begin{enumerate}[(1)]
  \item $A_P$ is a valuation ring for every maximal ideal $P$ of $A$;
  \item an $A$-module is flat if and only if it is torsion-free. (An integral
        domain satisfying these conditions is called a \emph{Pr\"ufer domain}.)
\end{enumerate}
\end{exercise}

\begin{exercise}\label{ex:11.10-ex}
A finite torsion-free module over a Dedekind ring is projective, and is
isomorphic to a direct sum of ideals.
\end{exercise}

% ============ §12 Krull rings (原书 pp.86-91) ============
\section{Krull rings}\label{sec:12}\index{Krull rings, section 12}

Let $A$ be an integral domain and $K$ its field of fractions. We write
$K^*$ for the multiplicative group of $K$. We say that $A$ is a \emph{Krull ring}\index{Krull ring} if there is
a family $\mathcal{F} = \{R_\lambda\}_{\lambda \in \Lambda}$ of DVRs of $K$ such that the following two conditions
hold, where we write $v_\lambda$ for the normalised additive valuation correspond-
ing to $R_\lambda$:
\begin{enumerate}[(1)]
  \item $A = \bigcap_\lambda R_\lambda$;
  \item for every $x \in K^*$ there are at most a finite number of $\lambda \in \Lambda$ such that
        $v_\lambda(x) \neq 0$.
\end{enumerate}
The family $\mathcal{F}$ of DVRs is said to be a \emph{defining family} of $A$. Since DVRs are
completely integrally closed (see Ex.~\ref{ex:9.5-ex}), so are Krull rings. If $A$ is a Krull
ring then for any subfield $K' \subset K$ the intersection $A \cap K'$ is again Krull.

\begin{theorem}\label{thm:12.1}
If $A$ is a Krull ring and $S \subset A$ a multiplicative set, then $A_S$ is
again Krull. If $\mathcal{F} = \{R_\lambda\}_{\lambda \in \Lambda}$ is a defining family of $A$ then the subfamily
$\{R_\lambda\}_{\lambda \in \Gamma}$, where $\Gamma = \{\lambda \in \Lambda \mid R_\lambda \supset A_S\}$, is a defining family of $A_S$.
\end{theorem}

\begin{proof}
Setting $m_\lambda$ for the maximal ideal of $R_\lambda$, we have
\[
  \bigcap_{\lambda \in \Gamma} m_\lambda = 0.
\]
Let $0 \neq x \in \bigcap_{\lambda \in \Gamma} R_\lambda$; there are at most finitely many $\lambda \in \Lambda$ such that $v_\lambda(x) < 0$;
let $\Delta = \{\lambda_1, \ldots, \lambda_n\}$ be the set of these. If $s \in S$ then $s \notin m_\lambda$, hence we can find
$t_\lambda \in m_\lambda \cap S$. Replacing $t_\lambda$ by a suitable power, we can assume that $v_\lambda(t_\lambda x) \ge 0$.
We then set $t = \prod_{\lambda \in \Delta} t_\lambda$, so that for every $\lambda \in \Lambda$ we have $v_\lambda(tx) \ge 0$, and
therefore $tx \in A$; but on the other hand $t \in S$ so that $x \in A_S$, and we have
proved that $A_S \supset \bigcap_{\lambda \in \Gamma} R_\lambda$. The opposite inclusion is obvious. The finiteness
condition (2) holds for $A$, so also for the subset $\Gamma$.
\end{proof}

Krull rings defined by a finite number of DVRs have a simple structure.

\begin{lemma}[Nagata]\label{lem:12.1}
Let $K$ be a field and $R_1, \ldots, R_n$ valuation rings of $K$;
set $A = \bigcap R_i$. Then for any given $a \in K$ there exists a natural number $s \ge 2$
such that
\[
  (1 + a + \cdots + a^{s-1})^{-1} \quad \text{and} \quad a \cdot (1 + a + \cdots + a^{s-1})^{-1}
\]
both belong to $A$.
\end{lemma}

\begin{proof}
We consider separately each $R_i$. Note first that $(1 - a)(1 + a + \cdots
+ a^{s-1}) = 1 - a^s$. If $a \notin R_i$ then $a^{-1} \in m_i$, and any $s \ge 2$ will do. If $a \in R_i$ then
provided that there does not exist $t$ such that $1 - a^t \in m_i$, any $s \ge 2$ will do.
If $1 - a \in m_i$ then any $s$ which is not a multiple of the characteristic of $R_i/m_i$
will do. If on the other hand $1 - a \notin m_i$ but $1 - a^t \in m_i$ for some $t > 2$, letting $t_0$
be the smallest value of $t$ for which this happens, we see that $1 - a^s \in m_i$
only for $s$ multiples of $t_0$, so that we only have to avoid these. Thus for each $i$
the bad values of $s$ (if any) are multiples of some number $d_i > 1$, so that
choosing $s$ not divisible by any of these $d_i$ we get the result.
\end{proof}

\begin{theorem}\label{thm:12.2}
Let $K$ be a field and $R_1, \ldots, R_n$ valuation rings of $K$ such
that $R_i \nsubseteq R_j$ for $i \neq j$; set $m_i = \rad(R_i)$. Then the intersection $A = \bigcap_{i=1}^{n} R_i$ is
a semilocal ring, having $p_i = m_i \cap A$ for $i = 1, \ldots, n$ as its only maximal
ideals; moreover $A_{p_i} = R_i$. If each $R_i$ is a DVR then $A$ is a principal ideal
ring.
\end{theorem}

\begin{proof}
The inclusion $A_{p_i} \subset R_i$ is obvious. For the opposite inclusion, let
$a \in R_i$; choosing $s \ge 2$ as in the lemma, and setting $u = (1 + a + \cdots + a^{s-1})^{-1}$
we get $u \in A$ and $au \in A$. Obviously $u$ is a unit of $R_i$, so that $u \in A - p_i$
and $a = (au)/u \in A_{p_i}$. This proves that $A_{p_i} = R_i$. It follows from this
that there are no inclusions among $p_1, \ldots, p_n$. If $I$ is an ideal of $A$ not
contained in any $p_i$ then (by Ex.~\ref{ex:1.6-ex}) there exists $x \in I$ not contained in
$\bigcup_{i=1}^{n} p_i$; then $x$ is a unit in each $R_i$, and hence in $A$, so that $I = A$. Thus
$p_1, \ldots, p_n$ are all the maximal ideals of $A$.

If each $R_i$ is a DVR then we have $m_i \neq m_i^{(2)}$, and hence $p_i \neq p_i^{(2)}$ (where $p^{(2)}$
denotes $p^2A_p \cap A$). Thus there exists $x_i \in p_i$ such that $x_i \notin p_i^{(2)}$, and $x_i \notin p_j$ for
$i \neq j$; then $p_i = x_iA$. If $I$ is any ideal of $A$ and $IR_i = x_i^{e_i}R_i$ for $i = 1, \ldots, n$ then it
is easy to see that $I = x_1^{e_1} \cdots x_n^{e_n}A$.
\end{proof}

If a Krull ring $A$ is defined by an infinite number of DVRs then the
defining family of DVRs is not necessarily unique, but the following
theorem tells us that among them there is a minimal family.

\begin{theorem}\label{thm:12.3}
Let $A$ be a Krull ring, $K$ its field of fractions, and $p$ a height $1$
prime ideal of $A$; then if $\mathcal{F} = \{R_\lambda\}_{\lambda \in \Lambda}$ is a family of DVRs of $K$ defining $A$, we
must have $A_p \in \mathcal{F}$. If we set $\mathcal{F}_0 = \{A_p \mid p \in \Spec A \text{ and } \htop p = 1\}$ then $\mathcal{F}_0$ is a
defining family of $A$. Thus $\mathcal{F}_0$ is the minimal defining family of DVRs of $A$.
\end{theorem}

\begin{proof}
By Theorem 1, $A_p$ is a Krull ring defined by the subfamily $\mathcal{F}_1 =
\{R_\lambda \mid A_p \subset R_\lambda\} \subset \mathcal{F}$; if $A_p \subset R_\lambda$ then the elements of $A - p$ are units of $R_\lambda$,
so that $p \supset m_\lambda \cap A$. If $m_\lambda \cap A = (0)$ then $R_\lambda \supset K$ which is a contradiction,
hence $m_\lambda \cap A \neq (0)$; since $\htop p = 1$, we must have $p = m_\lambda \cap A$. Thus if we fix
some $0 \neq x \in p$, then $v_\lambda(x) > 0$ for all $R_\lambda \in \mathcal{F}_1$, and hence $\mathcal{F}_1$ is a finite set,
Now by the previous theorem and Ex.~\ref{ex:10.5-ex}, the elements of $\mathcal{F}_1$ correspond
bijectively with maximal ideals of $A_p$, and $\mathcal{F}_1$ has just one element $A_p$. Thus
$A_p \in \mathcal{F}$, in other words $\mathcal{F}_0 \subset \mathcal{F}$.

To prove that $\mathcal{F}_0$ is a defining family of DVRs of $A$ it is enough to show that
$A \supset \bigcap_{\htop p = 1} A_p$. That is, it is enough to prove the implication
\[
  \text{for } a, b \in A \text{ with } a \neq 0, \quad b \in aA_p \text{ for all } A_p \in \mathcal{F}_0 \Rightarrow b \in aA.
\]
As one sees easily, this is equivalent to saying that $aA$ can be written as
the intersection of height $1$ primary ideals. The set of $R_\lambda \in \mathcal{F}$ such that
$aR_\lambda \neq R_\lambda$ is finite, so we write $R_1, \ldots, R_r$ for this. If we set
\[
  aR_i \cap A = q_i \quad \text{and} \quad \rad(R_i) \cap A = p_i
\]
then $q_i$ is a primary ideal belonging to $p_i$ for each $i$, and $aA = q_1 \cap \cdots \cap q_r$.
Eliminating redundant terms from this expression, we get an irredundant
expression, say $aA = q_1 \cap \cdots \cap q_r$. It is enough to show that then $\htop p_i = 1$
for $1 \le i \le r$. By contradiction, suppose that $\htop p_1 > 1$. By Theorem 1, $A_{p_1}$ is
a Krull ring with defining family $\mathcal{F}' = \{R_\lambda \in \mathcal{F} \mid A_{p_1} \subset R_\lambda\}$, but is not itself
a DVR, and hence by Theorem 2, $\mathcal{F}'$ is infinite. Thus there exists $R' \in \mathcal{F}'$
such that $aR' = R'$; we set $p' = \rad(R') \cap A$. We have $a \notin p'$, and $A_{p'} \subset R'$
implies that $p' \subset p_1$. Now by assumption $aA = q_1 \cap \cdots \cap q_r$ is irredundant, and $R_1$ is
a DVR, so that $(\rad(R_1))^\nu \subset aR_1$ for some $\nu > 0$, and hence $p_1^\nu \subset q_1$.
Therefore there exists an $i \ge 0$ such that
\[
  aA \neq p'^i \cap q_1 \cap \cdots \cap q_r \quad \text{and} \quad aA \supset p'^{i+1} \cap q_1 \cap \cdots \cap q_r.
\]
Hence there exist $b \in A$ such that $b \notin aA$ but $bp'^i \subset aA$. In particular $bp' \subset aA$,
but since $a$ is a unit of $R'$ we have
\[
  (b/a)p' \subset A \cap \rad(R') = p'.
\]
Taking $0 \neq c \in p'$ then for every $n > 0$ we have $(b/a)^nc \in p' \subset A$, and since $A$ is
completely integrally closed, $b/a \in A$. This is a contradiction, and it proves
that $\htop p_i = 1$ for $1 \le i \le r$.
\end{proof}

\begin{corollary*}
Let $A$ be a Krull ring and $\mathcal{P}$ the set of height $1$ prime ideals of $A$.
For $0 \neq a \in A$ set $v_p(a) = n_p$; then
\[
  aA = \bigcap_{p \in \mathcal{P}} p^{(n_p)},
\]
where $p^{(n)}$ denotes the symbolic $n$th power $p^nA_p \cap A$.
\end{corollary*}

\begin{proof}
According to the theorem we have $aA = \bigcap_p (aA_p \cap A)$,
but $aA_p = p^{n_p}A_p$, so that $aA_p \cap A = p^{(n_p)}$.
\end{proof}

\begin{theorem}\label{thm:12.4}
\begin{enumerate}[(i)]
  \item A Noetherian normal domain is a Krull ring.
  \item Let $A$ be an integral domain, $K$ its field of fractions, and $L$ an extension
        field of $K$. If $(A_i)_{i \in I}$ is a family of Krull rings contained in $L$ satisfying
        the two conditions (1) $A = \bigcap A_i$ and (2) given any $0 \neq a \in A$ we have
        $aA_i = A_i$ for all but finitely many $i$, then $A$ is a Krull ring.
  \item If $A$ is a Krull ring then so is $A[X]$ and $A[[X]]$.
\end{enumerate}
\end{theorem}

\begin{proof}
(i) This follows from Theorem~\ref{thm:11.5} and the fact that for any non-
zero $a \in A$ there are only finitely many height $1$ prime ideals containing $aA$
(because these are the prime divisors of $aA$).

(ii) is easy, and we leave it to the reader.

(iii) $K[X]$ is a principal ideal ring and therefore a Krull ring. Moreover,
if we let $\mathcal{P}$ be the set of height $1$ prime ideals of $A$ then for $p \in \mathcal{P}$
the ideal $p[X]$ is prime in $A[X]$, and by Theorem~\ref{thm:11.2}, (3), the local ring
$A[X]_{p[X]}$ is a DVR of $K(X)$. (If we write $v$ for the additive valuation of $K$
corresponding to the valuation ring $A_p$, we can extend $v$ to an additive
valuation of $K(X)$ by setting $v(F(X)) = \min\{v(a_i)\}$ for a polynomial
$F(X) = a_0 + a_1X + \cdots + a_rX^r$ (with $a_i \in K$),
and $v(F/G) = v(F) - v(G)$ for a rational function $F(X)/G(X)$; then the
valuation ring of $v$ in $K(X)$ is $A[X]_{p[X]}$.) Now we have
\[
  K[X] \cap A[X]_{p[X]} = A_p[X]
\]
(prove this!), and so
\[
  A[X] = \bigcap_{p \in \mathcal{P}} A_p[X],
\]
by (ii) this is a Krull ring.

Now for $A[[X]]$, let $\{R_\lambda\}_{\lambda \in \Lambda}$ be a family of DVRs of $K$ defining $A$; then
inside $K[[X]]$ we have $A[[X]] = \bigcap_\lambda R_\lambda[[X]]$; also by Ex.~\ref{ex:9.5-ex}, $R_\lambda[[X]]$ is an
integrally closed Noetherian ring, and is therefore a Krull ring by (i).
However, we cannot use (ii) as it stands, since $X$ is a non-unit of all the rings
$R_\lambda[[X]]$, so we set $R_\lambda[[X]][X^{-1}] = B_\lambda$ and note that $A[[X]] = K[[X]] \cap
(\bigcap_\lambda B_\lambda)$; now the hypothesis in (ii) is easily verified. Indeed,
\[
  F(X) = a_rX^r + a_{r+1}X^{r+1} + \cdots \in A[[X]] \quad \text{with } a_r \neq 0
\]
is a non-unit of $B_\lambda$ if and only if $a_r$ is a non-unit of $R_\lambda$, and there are only
finitely many such $\lambda$. Therefore $A[[X]]$ is a Krull ring.
\end{proof}

\begin{remark}
Note that the field of fractions of $A[[X]]$ is in general smaller
than the field of fractions of $K[[X]]$.
\end{remark}

\begin{remark}
The $B_\lambda$ occurring above are Euclidean rings ([B7], \S1, Ex. 9).
\end{remark}

\begin{theorem}\label{thm:12.5}
The notions of Dedekind ring and one-dimensional Krull
ring coincide.
\end{theorem}

\begin{proof}
A Dedekind ring is a normal Noetherian domain, and therefore a
Krull ring. Conversely, if $A$ is a one-dimensional Krull ring, let us prove that
$A$ is Noetherian. Let $I$ be a non-zero ideal of $A$, and let $0 \neq a \in I$. If we can
prove that $A/aA$ is Noetherian then $I/aA$ is finitely generated, and thus so is
$I$. By the corollary of Theorem 3 we can write $aA = q_1 \cap \cdots \cap q_r$, where
$q_i$ are symbolic powers of prime ideals $p_i$ and $p_i \neq p_j$ if $i \neq j$; now since
$\dim A = 1$ each $p_i$ is maximal and we have
\[
  A/aA = A/q_1 \times \cdots \times A/q_r
\]
by Theorem~\ref{thm:1.3} and Theorem~\ref{thm:1.4}. But $A/q_i$ is a local ring with maximal
ideal $p_i/q_i$, and hence $A/q_i \cong A_{p_i}/q_iA_{p_i}$; now since each $A_{p_i}$ is a DVR, $A/aA$ is
Noetherian (in fact even Artinian). Hence $A$ is one-dimensional Noetherian
integral domain, and is normal, and is therefore a Dedekind ring.
\end{proof}

\begin{theorem}\label{thm:12.6}
Let $A$ be a Krull ring, $K$ its field of fractions, and write $\mathcal{P}$
for the set of height $1$ prime ideals of $A$. Suppose given any $p_1, \ldots, p_r \in \mathcal{P}$
and $e_1, \ldots, e_r \in \Z$. Then there exists $x \in K$ satisfying
\[
  v_{p_i}(x) = e_i \quad \text{for } 1 \le i \le r
\]
and
\[
  v_p(x) \ge 0 \quad \text{for all } p \in \mathcal{P} - \{p_1, \ldots, p_r\}.
\]
Here $v_{p_i}$ and $v_p$ stand for the normalised additive valuations of $K$
corresponding to $p_i$ and $p$.
\end{theorem}

\begin{proof}
If $y_i \in A$ is chosen so that $y_i \in p_i$ but $y_i \notin p_i^{(2)}$ and $y_i \notin p_j$ for
$j \neq i$, then $v_{p_j}(y_i) = \delta_{ij}$ for $1 \le i, j \le r$. Similarly we choose $y_1, \ldots, y_r \in A$ such
that $v_{p_i}(y_j) = \delta_{ij}$. Then we set
\[
  y = \prod_{i=1}^{r} y_i^{e_i};
\]
let $p_1', \ldots, p_s'$ be all the primes $p \in \mathcal{P} - \{p_1, \ldots, p_r\}$ for which $v_p(y) \neq 0$.
Then choosing for each $j = 1, \ldots, s$ an element $t_j \in p_j'$ not belonging to
$p_1 \cup \cdots \cup p_r$, and taking $\nu$ to be sufficiently large, we see that
\[
  x = y(t_1 \cdots t_s)^\nu
\]
satisfies the requirements of the theorem.
\end{proof}

\begin{theorem}[Y. Mori and J. Nishimura]\label{thm:12.7}
Let $A$ and $\mathcal{P}$ be as in the previous
theorem. If $A/p$ is Noetherian for every $p \in \mathcal{P}$ then $A$ is Noetherian.
\end{theorem}

\begin{proof}[Proof (J. Nishimura)]
As in the proof of Theorem 5, it is enough to show
that $A/p^{(n)}$ is Noetherian for $p \in \mathcal{P}$ and any $n > 0$. If $n = 1$ this holds by
hypothesis. For $n > 1$ we proceed as follows. Applying the previous
theorem with $r = 1$ and $e = -1$ we can find an element $x$ in the field of
fractions of $A$ such that $v_p(x) = 1$ and $v_q(x) < 0$ for every other $q \in \mathcal{P}$. Set
$B = A[x]$. If $y \in p$ then $y/x \in A$ so that $y \in xB$, and conversely $B \subset A_p$, and
$xB \subset pA_p$, so that $p = xB \cap A$. Moreover, $B = A + xB$ and so
\[
  B/xB = A/p.
\]
Now $x^iB/x^{i+1}B \cong B/xB$ for each $i$, so that by induction on $i$ we see that
$B/x^iB$ is a Noetherian $B$-module for each $i$, and hence a Noetherian ring.
Now we have
\[
  A/(x^nB \cap A) \hookrightarrow B/x^nB,
\]
and $B/x^nB$ is finite over $A/(x^nB \cap A)$, being generated by the images of
$1, x, \ldots, x^{n-1}$, so that by the Eakin--Nagata theorem (Theorem~\ref{thm:3.7}),
$A/(x^nB \cap A)$ is a Noetherian ring; therefore its quotient $A/p^{(n)}$ is also
Noetherian.
\end{proof}

\begin{remark}
If $A$ is a Noetherian integral domain and $K$ its field of fractions,
then the integral closure of $A$ in $K$ is a possibly non-Noetherian Krull ring
([N1], (33.10)). This was proved by Y. Mori (1952) in the local case, and
in the general case by M. Nagata (1955). Theorem~\ref{thm:12.7} was proved by
Mori [1] in 1955 as a theorem on the integral closure of Noetherian
rings. His proof was correct (in spite of a number of easily rectifiable
inaccuracies), and was an extremely interesting piece of work, but due to
its difficulty, and the fact that it appeared in an inaccessible journal, the
result was practically forgotten. After Marot [1], (1973) applied it
successfully, Mori's work attracted attention once more, and J. Nishimura
[1], (1975) reformulated the result as above as a theorem on Krull rings
and gave an elegant proof.

More results on Krull rings can be found in [N1], [B7], [F], among
others.
\end{remark}

% ============ 习题 ============
\section*{Exercises to \S\ref{sec:12}}
\emph{Prove the following propositions.}

\begin{exercise}\label{ex:12.1-ex}
Let $K \subset L$ be a finite extension of fields, and $R$ a valuation ring of $K$. Then
there are a finite number of valuation rings of $L$ dominating $R$, and if $L$ is a
normal extension of $K$ then these are all conjugate to one another under
elements of the Galois group $\Aut_K(L)$.
\end{exercise}

\begin{exercise}\label{ex:12.2-ex}
Let $R$ be a valuation ring of a field $K$, and let $K \subset L$ be a (possibly infinite)
algebraic extension; write $\bar{R}$ for the integral closure of $R$ in $L$. Then the
localisation of $\bar{R}$ at a maximal ideal is a valuation ring dominating $R$, and
conversely every valuation ring of $L$ dominating $R$ is obtained in this way.
\end{exercise}

\begin{exercise}\label{ex:12.3-ex}
Let $A$ be a Krull ring, $K$ its field of fractions, and $K \subset L$ a finite extension
field; if $B$ is the integral closure of $A$ in $L$, then $B$ is also a Krull ring.
\end{exercise}

\begin{exercise}\label{ex:12.4-ex}
Let $A$ be an integral domain and $K$ its field of fractions. For $I$ a fractional
ideal, write $\tilde{I} = (I^{-1})^{-1}$. If $I = \tilde{I}$ we say that $I$ is \emph{divisorial}. If $A$ is a Krull
ring, then an ideal of $A$ is divisorial if and only if it can be expressed as the
intersection of a finite number of height $1$ primary ideals.
\end{exercise}

